\topic{本章实验、故意破坏与练习}

\begin{experimentbox}{从符号追到 ROM byte}
选择复位入口符号，依次记录目标文件地址、链接地址、ROM 文件偏移和最终三个
byte；再运行 QEMU，确认首条串口日志带时间出现。
\end{experimentbox}

\begin{faultinjectionbox}{保持 DLAB=1 就发送字符}
DLAB 为 1 时，base 端口访问除数低 byte，不是发送保持寄存器。预期字符不出现
且波特率配置被改坏。清 DLAB 后再等 THRE 并发送。
\end{faultinjectionbox}

\begin{pitfallbox}
源码符号不等于最终 byte；THRE 要有有界轮询；串口时间戳不能早于可用时钟来源。
\end{pitfallbox}

\textbf{练习}
\begin{enumerate}[leftmargin=2.2em]
  \item 128 KiB ROM 的十六进制大小是多少？
  \item THRE 未就绪时应继续写还是等待？
  \item 为什么超时日志要打印最后一次 LSR？
\end{enumerate}
\begin{answerbox}
1. \texttt{0x20000}；2. 有界等待；3. 区分设备一直 busy、读错端口和线路状态
异常。
\end{answerbox}
